% БҮЛЭГ 2: СИСТЕМИЙН ШААРДЛАГА БА ШИНЖИЛГЭЭ
\chapter{СИСТЕМИЙН ШААРДЛАГА БА ШИНЖИЛГЭЭ}
\label{Chapter2}
\pagecolor{white}

%===============================================================
\section{2.1 Системийн үйл ажиллагааны тухай}

Санал болгож буй систем нь дараах бүрэн урсгалыг хэрэгжүүлнэ:

\begin{enumerate}
    \item Үйлчлүүлэгч ресторанд орж ширээнд суухад ширээн дэлгэц дээр QR код байрлана;
    \item Хэрэглэгч смартфоны камераар QR кодыг уншуулна;
    \item Мобайл аппликейшн автоматаар нээгдэж тухайн рестораны дижитал цэс 3 секундын
          дотор харагдана;
    \item Хэрэглэгч хоол сонгож сагсанд нэмэн ``Захиалга өгөх'' товчийг дарна;
    \item Захиалга нь WebSocket протоколоор ажилтны дэлгэцэнд бодит цаг хугацаанд дамжина;
    \item Цайчин захиалгыг баталгаажуулахад зочны дэлгэцэнд мэдэгдэл ирнэ;
    \item Тогооч хоол бэлтгэж ``Бэлэн боллоо'' тэмдэглэхэд зочны дэлгэцэнд мэдэгдэл ирнэ;
    \item Зочин тооцоо хийлгэхэд кассир тооцоог баталгаажуулна.
\end{enumerate}

%===============================================================
\section{2.2 Системийг ашиглах хэрэглэгчид}

\begin{table}[htbp]
\caption{Системийн хэрэглэгчдийн үүргийн тодорхойлолт}
\label{tab:roles}
\centering
\begin{tabularx}{\textwidth}{|>{\bfseries}X|X|X|}
\hline
\rowcolor{TableHeader}\textcolor{white}{\textbf{Хэрэглэгч}}
  & \textcolor{white}{\textbf{Хандах нөөц}}
  & \textcolor{white}{\textbf{Боломжит үйлдэл}} \\ \hline
Зочин (Guest)       & Цэс, захиалга        & Харах, захиалга үүсгэх \\ \hline
Цайчин (Waiter)     & Захиалга, ширээ      & Харах, баталгаажуулах, шинэчлэх \\ \hline
Тогооч (Chef)       & Захиалга             & Харах, бэлтгэлт шинэчлэх \\ \hline
Кассир (Cashier)    & Тооцоо, захиалга     & Харах, төлбөр гүйцэтгэх \\ \hline
Менежер (Manager)   & Бүх нөөц            & Харах, нэмэх, засварлах \\ \hline
Администратор       & Бүх нөөц            & Бүрэн хяналт \\ \hline
\end{tabularx}
\end{table}

%===============================================================
\section{2.3 Функционал шаардлага}

\textbf{ФШ-01: QR Кодоор Цэс Нээх}\\
Хэрэглэгч ширээн дэлгэц дээрх QR кодыг смартфоны камераар уншуулна.
3 секундын дотор тухайн ширээний мэдээлэл болон ресторан цэс нэн даруй нээгдэнэ.
Цэс нь ангиллаар бүлэглэгдсэн, хоол бүрийн зураг, нэр, тайлбар, үнэ агуулсан байна.
QR кодын хүчинтэй хугацаа 24 цаг байх ба хугацаа дуусвал ажилтан шинэчилнэ.

\textbf{ФШ-02: Захиалга Үүсгэх}\\
Хэрэглэгч цэснээс хоол сонгон сагсанд нэмэх, хасах боломжтой байна.
Хоол бүрд тусгай хүсэлт бичих боломжтой байна.
``Захиалга өгөх'' товч дарахад нийт дүн болон захиалгын мэдээллийн баталгаажуулалтын
дэлгэц харагдана.

\textbf{ФШ-03: Бодит Цагийн Мэдэгдэл (Зочин)}\\
Захиалга баталгаажсан, бэлтгэж байна, бэлэн боллоо, хүргэгдсэн зэрэг статус
өөрчлөлт бүрийг зочны дэлгэцэнд шуурхай харуулна. Дэлгэц унтарсан үед Push
Notification илгээнэ.

\textbf{ФШ-04: Бодит Цагийн Мэдэгдэл (Ажилтан)}\\
Шинэ захиалга ирэхэд ажилтны дэлгэцэнд дохиолол (дуу + чичиргэх) болон мэдэгдэл
харагдана. Захиалга нь ширээний нэр, хоолны жагсаалт, тусгай хүсэлт, нийт дүн
агуулсан байна.

\textbf{ФШ-05: Захиалгын Статус Удирдлага}\\
Захиалгын статусын шилжилт нь зөвхөн тодорхой дарааллаар явагдана:
pending $\rightarrow$ confirmed $\rightarrow$ preparing $\rightarrow$ ready $\rightarrow$ served $\rightarrow$ paid.
Статус өөрчлөлт бүр WebSocket-р зочны дэлгэцэнд шуурхай мэдэгдэнэ.

\textbf{ФШ-06: Цэс Удирдлага}\\
Менежер цэсийн ангилал болон хоолыг нэмэх, засварлах, устгах боломжтой байна.
Хоолны байдал (байна/байхгүй) бодит цаг хугацаанд шинэчлэх боломжтой байна.
Шинэчлэлт нэн даруй хэрэглэгчдийн дэлгэцэнд тусна.

%===============================================================
\section{2.4 Функционал бус шаардлага}

\subsection{2.4.1 Програм хангамжийн чанарын шаардлага}

\begin{table}[htbp]
\caption{Гүйцэтгэлийн шаардлагууд}
\label{tab:perf}
\centering
\begin{tabularx}{\textwidth}{|X|X|X|}
\hline
\rowcolor{TableHeader}\textcolor{white}{\textbf{Метрик}}
  & \textcolor{white}{\textbf{Зорилт}}
  & \textcolor{white}{\textbf{Хэмжих аргачлал}} \\ \hline
API хариу хугацаа    & Дундаж 200 ms, 95th\% 400 ms & Apache JMeter \\ \hline
WebSocket хоцролт    & Дундаж 50 ms доор             & Серверийн лог \\ \hline
Зэрэг хэрэглэгч      & Нэг рестораны 80 хэрэглэгч   & Load test \\ \hline
Uptime               & 99.5\%                        & Monitoring \\ \hline
QR нээх хугацаа      & 3 секунд доор                 & Хэрэглэгчийн тест \\ \hline
\end{tabularx}
\end{table}

\textbf{Аюулгүй байдал:} Бүх дамжуулалт HTTPS (TLS 1.3) протоколоор шифрлэгдэнэ.
Нэвтрэх оролдлогыг хязгаарлах rate limiting, SQL injection хамгаалалт, input
validation зэрэг хамгаалалтын механизмуудыг хэрэгжүүлнэ \cite{owasp2023}.

\textbf{Платформын нийцтэй байдал:} iOS 14.0 болон дээш, Android 10 (API 29) болон
дээш дэмжинэ.

%===============================================================
\section{2.5 Ашигласан технологи, техникийн шаардлага}

\subsection{2.5.1 Технологийн шаардлага}

\textbf{Програм ажиллах үйлдлийн системийн сонголт}

API сервер нь Ubuntu 22.04 LTS үйлдлийн систем дээр ажиллана. Ubuntu нь серверийн
орчинд найдвартай байдал, аюулгүй байдал болон олон нийтийн дэмжлэгийн хувьд дэлхий
дээр хамгийн өргөн хэрэглэгддэг нэг юм.

\textbf{Програм хөгжүүлэх хэлний сонголт}

TypeScript нь JavaScript дээр суурилсан type-safe програмчлалын хэл бөгөөд мобайл
болон сервер аль алины кодыг нэг хэлд хөгжүүлэх боломжтой. Type inference нь
компайлын үед алдааг илрүүлэх давуу талтай.

\textbf{Өгөгдлийн сангийн удирдах системийн сонголт}

PostgreSQL нь ACID баталгаа, массив төрлийн өгөгдөл (алергийн жагсаалт), JSON
дэмжлэг зэрэг онцлогоороо энэхүү системд хамгийн тохиромжтой.

\subsection{2.5.2 Техникийн шаардлага}

\begin{table}[htbp]
\caption{Серверийн техникийн шаардлага}
\label{tab:hardware}
\centering
\begin{tabularx}{0.85\textwidth}{|X|X|}
\hline
\rowcolor{TableHeader}\textcolor{white}{\textbf{Бүрэлдэхүүн хэсэг}}
  & \textcolor{white}{\textbf{Шаардлага}} \\ \hline
CPU          & 2 virtual core, 2.4 GHz+     \\ \hline
RAM          & 2 GB+                         \\ \hline
Хадгалалт   & 20 GB SSD                     \\ \hline
Сүлжээ      & 10 Mbps тогтвортой            \\ \hline
Үйлдлийн систем & Ubuntu 22.04 LTS           \\ \hline
\end{tabularx}
\end{table}

\textbf{Оролт гаралтын төхөөрөмжийн шаардлага}

Мобайл клиент нь iOS 14.0+, Android 10+ дэмжих бөгөөд хамгийн багадаа 2 GB RAM,
QR уншигч нийцтэй камертай байх шаардлагатай. QR код ширээн дэлгэцэнд 5$\times$5 см
хэмжээнд хэвлэж, ламинат хийж байрлуулна.

%===============================================================
\section{2.6 Юзкейс диаграм}

Системийн юзкейс диаграмыг Зураг~\ref{fig:usecase} дээр харуулав.

\begin{figure}[htbp]
\centering
\begin{tikzpicture}[
    actor/.style={rectangle, draw, minimum width=1.2cm, minimum height=0.6cm, align=center, font=\small},
    usecase/.style={ellipse, draw, minimum width=2.8cm, minimum height=0.8cm, align=center, font=\small},
    system/.style={rectangle, draw, dashed, inner sep=10pt}
]

% Actors
\node[actor] (guest) at (-5, 2) {Зочин};
\node[actor] (waiter) at (-5, -1) {Цайчин};
\node[actor] (chef) at (7, 0.5) {Тогооч};
\node[actor] (manager) at (7, -2) {Менежер};

% System boundary
\node[system, minimum width=9cm, minimum height=7cm] (sys) at (1, 0) {};
\node[above] at (sys.north) {\textbf{Рестораны захиалгын систем}};

% Use cases
\node[usecase] (qr) at (0, 3) {QR уншуулах};
\node[usecase] (menu) at (0, 1.8) {Цэс харах};
\node[usecase] (order) at (0, 0.5) {Захиалга өгөх};
\node[usecase] (status) at (0, -0.8) {Статус харах};
\node[usecase] (confirm) at (3, 1.5) {Захиалга баталгаажуулах};
\node[usecase] (kitchen) at (3, 0) {Хоол бэлтгэх};
\node[usecase] (menuMgr) at (3, -2) {Цэс удирдах};

% Connections
\draw (guest) -- (qr);
\draw (guest) -- (menu);
\draw (guest) -- (order);
\draw (guest) -- (status);
\draw (waiter) -- (confirm);
\draw (waiter) -- (status);
\draw (chef) -- (kitchen);
\draw (manager) -- (menuMgr);

\end{tikzpicture}
\caption{Системийн юзкейс диаграм}
\label{fig:usecase}
\end{figure}

%===============================================================
\section{2.7 Юзкейс диаграмын тодорхойлолт}

\begin{table}[htbp]
\caption{Голлох юзкейсүүдийн тодорхойлолт}
\label{tab:usecasedesc}
\centering
\begin{tabularx}{\textwidth}{|>{\bfseries}l|X|}
\hline
\rowcolor{TableHeader}\multicolumn{2}{|c|}{\textcolor{white}{\textbf{UC-01: QR Уншуулах}}} \\ \hline
Жүжигчин     & Зочин \\ \hline
Урьдчилсан нөхцөл & Ширээн дэлгэцэнд QR код байна, камерын зөвшөөрөл өгсөн \\ \hline
Амжилтын хувилбар & QR уншиж 3 секундын дотор цэс нээгдэнэ \\ \hline
Бүтэлгүйтлийн хувилбар & Token хугацаа дуусвал ``QR шинэчлүүлнэ'' мессеж гарна \\ \hline
\end{tabularx}
\end{table}

\begin{table}[htbp]
\centering
\begin{tabularx}{\textwidth}{|>{\bfseries}l|X|}
\hline
\rowcolor{TableHeader}\multicolumn{2}{|c|}{\textcolor{white}{\textbf{UC-02: Захиалга Өгөх}}} \\ \hline
Жүжигчин     & Зочин \\ \hline
Урьдчилсан нөхцөл & Цэс нээгдсэн, Session ID үүссэн \\ \hline
Амжилтын хувилбар & Захиалга ажилтны дэлгэцэнд хүрч дохиолол дуугарна \\ \hline
Бүтэлгүйтлийн хувилбар & Хоол байхгүй бол ``Одоогоор байхгүй'' мессеж гарна \\ \hline
\end{tabularx}
\end{table}

\begin{table}[htbp]
\centering
\begin{tabularx}{\textwidth}{|>{\bfseries}l|X|}
\hline
\rowcolor{TableHeader}\multicolumn{2}{|c|}{\textcolor{white}{\textbf{UC-03: Захиалга Баталгаажуулах}}} \\ \hline
Жүжигчин     & Цайчин \\ \hline
Урьдчилсан нөхцөл & Систем нэвтэрсэн, pending захиалга ирсэн \\ \hline
Амжилтын хувилбар & Захиалгыг confirmed болгоход зочид мэдэгдэл очно \\ \hline
Бүтэлгүйтлийн хувилбар & Зөвшөөрөлгүй статус шилжилт хийхийг систем хаана \\ \hline
\end{tabularx}
\end{table}
